%% titre: Méthodes — Montrer qu'un ensemble est fermé
%% matiere: Optimisation
%% niveau: Licence 3
%% annee: 2026-2027
%% ordre: 8
%% description: Méthodes pour démontrer qu'un ensemble est fermé à l'aide de la continuité, des images réciproques et des suites convergentes.

\section{Démarche}

\begin{methode}[Montrer qu'un ensemble $S\subset\R^n$ est fermé]
\begin{enumerate}[label=\textbf{Étape \arabic*.}]
    \item \textbf{Identifier les contraintes} : écrire l'ensemble sous la forme
    \[
    S=\{x\in\R^n:g_i(x)\leq 0,\ i=1,\ldots,m\}.
    \]

    \item \textbf{Vérifier la continuité} : montrer que les fonctions
    \[
    g_i:\R^n\to\R
    \]
    sont continues.

    \item \textbf{Utiliser l'image réciproque} : puisque $]-\infty,0]$ est fermé et que chaque $g_i$ est continue, on sait que
    \[
    g_i^{-1}(]-\infty,0])
    \]
    est fermé.

    \item \textbf{Conclure} : une intersection finie d'ensembles fermés étant fermée, on obtient
    \[
    \boxed{S\text{ est fermé}.}
    \]
\end{enumerate}
\end{methode}

\begin{attention}
\begin{itemize}
    \item L'image réciproque d'un fermé par une fonction continue est fermée.
    \item L'intersection finie d'ensembles fermés est fermée.
    \item Une réunion finie d'ensembles fermés est également fermée.
    \item Une inégalité stricte, telle que $g(x)<0$, définit en général un ensemble ouvert et non fermé.
    \item Un ensemble fermé n'est pas nécessairement borné.
\end{itemize}
\end{attention}

\section{Autre méthode : caractérisation séquentielle}

\begin{methode}[Utiliser les suites convergentes]
Un ensemble $S\subset\R^n$ est fermé si et seulement si toute suite $(x_k)$ d'éléments de $S$ qui converge dans $\R^n$ a sa limite dans $S$.

Pour appliquer cette caractérisation :
\begin{enumerate}[label=\textbf{Étape \arabic*.}]
    \item Considérer une suite arbitraire $(x_k)\subset S$ telle que
    \[
    x_k\longrightarrow x\in\R^n.
    \]
    \item Utiliser les contraintes vérifiées par chaque $x_k$.
    \item Passer à la limite en utilisant la continuité des fonctions de contraintes.
    \item Vérifier que la limite $x$ satisfait encore toutes les contraintes.
    \item Conclure que $x\in S$, donc que $S$ est fermé.
\end{enumerate}
\end{methode}

\section{Applications}

\begin{exemple}[Ensemble défini par des inégalités]
Considérons
\[
S=\left\{(x,y)\in\R^2:
\begin{array}{l}
x+y\leq 4,\\
x+3y\leq 9
\end{array}\right\}.
\]

Posons
\[
g_1(x,y)=x+y-4,\qquad
g_2(x,y)=x+3y-9.
\]

Les fonctions $g_1$ et $g_2$ sont affines, donc continues sur $\R^2$.

Or,
\[
S=g_1^{-1}(]-\infty,0])
\cap g_2^{-1}(]-\infty,0]).
\]

Puisque $]-\infty,0]$ est fermé, les deux images réciproques sont fermées. Leur intersection étant fermée, on obtient
\[
\boxed{S\text{ est fermé dans }\R^2.}
\]
\end{exemple}

\begin{exemple}[Ensemble défini par une égalité]
Considérons
\[
S=\{(x,y)\in\R^2:x^2+y^2=1\}.
\]

Posons
\[
g(x,y)=x^2+y^2-1.
\]

La fonction $g$ est continue et
\[
S=g^{-1}(\{0\}).
\]

Comme $\{0\}$ est un fermé de $\R$, son image réciproque par $g$ est fermée.

Ainsi,
\[
\boxed{S\text{ est fermé}.}
\]
\end{exemple}

\begin{attention}
Pour un ensemble défini par des contraintes mixtes
\[
S=\{x\in\R^n:g_i(x)\leq0,\ h_j(x)=0\},
\]
si toutes les fonctions $g_i$ et $h_j$ sont continues, alors $S$ est fermé.
\end{attention}
